\documentclass[a4paper,12pt]{article}

\usepackage[T2A]{fontenc}
\usepackage[utf8]{inputenc}
\usepackage[russian]{babel}

\usepackage{amsmath,amssymb,mathtools}
\usepackage{helvet}
\renewcommand{\familydefault}{\sfdefault}
\usepackage{icomma}
\usepackage{geometry}
\usepackage{microtype}
\usepackage{tikz}
\usetikzlibrary{arrows.meta,calc}
\usepackage{tabularx,booktabs}
\usepackage{enumitem}
\usepackage{hyperref}
\usepackage{parskip}
\usepackage{fancyhdr}
\usepackage{setspace}
\usepackage{xcolor}

\geometry{left=18mm,right=18mm,top=18mm,bottom=20mm}
\onehalfspacing

\definecolor{accent}{HTML}{008080}
\definecolor{darkaccent}{HTML}{004D4D}
\definecolor{textgray}{HTML}{333333}
\definecolor{softgray}{HTML}{F2F5F5}

\hypersetup{hidelinks}

\pagestyle{fancy}
\fancyhf{}
\fancyhead[L]{\textcolor{darkaccent}{Логарифмические неравенства}}
\fancyhead[R]{\textcolor{textgray}{ЕГЭ}}
\fancyfoot[C]{\thepage}

\setlist[itemize]{leftmargin=*,itemsep=0.25em,topsep=0.25em}
\setlist[enumerate]{leftmargin=*,itemsep=0.35em,topsep=0.35em}

\newcommand{\R}{\mathbb{R}}
\newcommand{\topic}{Логарифмические неравенства: модуль, ОДЗ и метод замены множителей}
\newcommand{\student}{Николь Саркисянц}
\newcommand{\teacher}{Лёвин Артём Александрович}

\begin{document}

\begin{titlepage}
\thispagestyle{empty}
\centering

\vspace*{35mm}

{\Large\textcolor{accent}{\textbf{Чек-лист}}\par}

\vspace{8mm}

{\Huge\textbf{\topic}\par}

\vspace{12mm}

{\large \teacher{} эксклюзивно для \student\par}

\vspace{8mm}

{\large \today\par}

\vfill

\begin{tikzpicture}[remember picture,overlay]
  \draw[
    line width=1.2pt,
    accent,
    opacity=0.65
  ]
  ([xshift=25mm,yshift=30mm]current page.south west)
  .. controls
  ([xshift=70mm,yshift=55mm]current page.south west)
  and
  ([xshift=-70mm,yshift=12mm]current page.south east)
  ..
  ([xshift=-25mm,yshift=38mm]current page.south east);
\end{tikzpicture}

\end{titlepage}

\section*{1. Главная идея темы}

Логарифмические неравенства повышенной сложности обычно опасны не самими логарифмами, а деталями:

\begin{itemize}
  \item областью допустимых значений;
  \item чётными степенями под логарифмом;
  \item появлением модуля;
  \item кратностью корней;
  \item выбором правильного равносильного перехода;
  \item аккуратным оформлением системы или совокупности.
\end{itemize}

\textbf{Главная цель:} превратить логарифмическое неравенство в обычное алгебраическое неравенство без потери ОДЗ.

\section*{2. Базовые ограничения для логарифма}

Для выражения

\[
\log_a f(x)
\]

обязательно:

\[
a>0,
\qquad
a\ne1,
\qquad
f(x)>0.
\]

Если основание является переменным, ограничения на основание нельзя пропускать.

Например:

\[
\log_x(2x-3)
\]

имеет ограничения:

\[
x>0,
\qquad
x\ne1,
\qquad
2x-3>0.
\]

То есть:

\[
x>\frac32.
\]

\section*{3. Чётная степень под логарифмом}

Если под логарифмом стоит чётная степень, её можно вынести перед логарифм, но основание степени превращается в модуль:

\[
\log_a\left(f(x)\right)^{2n}
=
2n\log_a|f(x)|.
\]

Например:

\[
\log_5(x-1)^2=2\log_5|x-1|.
\]

\textbf{Почему нужен модуль:} чётная степень делает выражение неотрицательным, но при вынесении степени нужно сохранить положительность результата.

\[
(x-1)^2=\bigl|x-1\bigr|^2.
\]

Поэтому:

\[
\log_5(x-1)^2
=
\log_5|x-1|^2
=
2\log_5|x-1|.
\]

\textbf{Важное ограничение:}

\[
(x-1)^2>0
\quad\Longleftrightarrow\quad
x\ne1.
\]

Точка \(x=1\) запрещена, потому что под логарифмом был бы ноль.

\section*{4. Пример с чётной степенью и модулем}

Решим неравенство:

\[
\log_5(x-1)^2\le2.
\]

Вынесем чётную степень:

\[
2\log_5|x-1|\le2.
\]

Сократим на \(2\):

\[
\log_5|x-1|\le1.
\]

Представим правую часть как логарифм:

\[
1=\log_5 5.
\]

Так как основание \(5>1\), знак неравенства сохраняется:

\[
|x-1|\le5.
\]

Решаем модульное неравенство:

\[
-5\le x-1\le5.
\]

\[
-4\le x\le6.
\]

Но нужно учесть ОДЗ:

\[
x\ne1.
\]

Ответ:

\[
\boxed{x\in[-4;1)\cup(1;6]}
\]

\section*{5. Модульные неравенства: быстрые переходы}

Модуль лучше не раскрывать по определению, если можно применить стандартный переход.

\subsection*{5.1. Модуль меньше числа}

Если

\[
|f(x)|\le g(x),
\]

то обычно переходим к системе:

\[
\begin{cases}
f(x)\le g(x),\\
f(x)\ge -g(x).
\end{cases}
\]

Если \(g(x)\) — положительное число, то это особенно удобно:

\[
|x-1|\le5
\quad\Longleftrightarrow\quad
-5\le x-1\le5.
\]

\subsection*{5.2. Модуль больше числа}

Если

\[
|f(x)|\ge g(x),
\]

то обычно переходим к совокупности:

\[
\left[
\begin{array}{l}
f(x)\ge g(x),\\
f(x)\le -g(x).
\end{array}
\right.
\]

Например:

\[
|x+2|\ge9
\quad\Longleftrightarrow\quad
\left[
\begin{array}{l}
x+2\ge9,\\
x+2\le-9.
\end{array}
\right.
\]

\[
\left[
\begin{array}{l}
x\ge7,\\
x\le-11.
\end{array}
\right.
\]

\[
x\in(-\infty;-11]\cup[7;+\infty).
\]

\section*{6. Кратность корня и знак на интервалах}

Если выражение содержит квадрат:

\[
(x-1)^2,
\]

то точка \(x=1\) является корнем чётной кратности.

\textbf{Правило:} при переходе через корень чётной кратности знак не меняется.

Например:

\[
(x-1)^2>0.
\]

Квадрат положителен при всех \(x\), кроме точки, где он равен нулю:

\[
x\in(-\infty;1)\cup(1;+\infty).
\]

\begin{center}
\begin{tikzpicture}[x=1.1cm,y=0.8cm]
  \draw[-{Stealth[length=2mm]}] (-4,0)--(4,0);
  \node[right] at (4,0) {$x$};
  \draw (0,0.12)--(0,-0.12);
  \node[below] at (0,-0.15) {$1$};

  \draw[very thick,accent] (-4,0.35)--(0,0.35);
  \draw[very thick,accent] (0,0.35)--(4,0.35);
  \draw[accent,fill=white,thick] (0,0.35) circle (2.3pt);

  \node[above,text=darkaccent] at (-2,0.45) {$+$};
  \node[above,text=darkaccent] at (2,0.45) {$+$};
\end{tikzpicture}
\end{center}

\section*{7. Неполное квадратное уравнение вместо дискриминанта}

Иногда не нужно раскрывать скобки и считать дискриминант.

Например:

\[
(x-1)^2\le25.
\]

Можно рассуждать через модуль:

\[
|x-1|\le5.
\]

\[
-5\le x-1\le5.
\]

\[
-4\le x\le6.
\]

Или через уравнение границ:

\[
(x-1)^2=25.
\]

\[
x-1=5
\quad\text{или}\quad
x-1=-5.
\]

\[
x=6
\quad\text{или}\quad
x=-4.
\]

Такой путь часто проще и надёжнее, чем дискриминант.

\section*{8. Метод замены множителей}

Метод замены множителей применяется к логарифмическим неравенствам с переменным основанием.

Если

\[
\log_a f(x)\;\vee\;\log_a g(x),
\]

где \(\vee\) — один из знаков \(>,<,\ge,\le\), то после записи ОДЗ можно заменить неравенство на:

\[
(a-1)\bigl(f(x)-g(x)\bigr)\;\vee\;0.
\]

Полная схема:

\[
\log_a f(x)\;\vee\;\log_a g(x)
\]

равносильно системе:

\[
\begin{cases}
a>0,\\
a\ne1,\\
f(x)>0,\\
g(x)>0,\\
(a-1)\bigl(f(x)-g(x)\bigr)\;\vee\;0.
\end{cases}
\]

\textbf{Главное преимущество:} логарифмы исчезают, а остаётся обычное алгебраическое неравенство.

\section*{9. Как подогнать неравенство под метод}

Метод замены множителей работает, когда слева и справа есть логарифмы с одинаковым основанием.

Если справа стоит число, его нужно представить в виде логарифма.

\subsection*{Пример с единицей}

\[
\log_x(2x-5)\ge1.
\]

Так как

\[
1=\log_x x,
\]

получаем:

\[
\log_x(2x-5)\ge\log_x x.
\]

Теперь применяем метод замены множителей:

\[
\begin{cases}
x>0,\\
x\ne1,\\
2x-5>0,\\
x>0,\\
(x-1)\bigl((2x-5)-x\bigr)\ge0.
\end{cases}
\]

Упрощаем последнее неравенство:

\[
(x-1)(x-5)\ge0.
\]

\section*{10. Если справа ноль}

Если справа стоит ноль, его тоже можно представить как логарифм:

\[
0=\log_a1.
\]

Например:

\[
\log_{x+1}(3x+2)<0.
\]

Представим:

\[
0=\log_{x+1}1.
\]

Тогда:

\[
\log_{x+1}(3x+2)<\log_{x+1}1.
\]

Применяем метод замены множителей:

\[
\begin{cases}
x+1>0,\\
x+1\ne1,\\
3x+2>0,\\
1>0,\\
(x+1-1)\bigl((3x+2)-1\bigr)<0.
\end{cases}
\]

\[
\begin{cases}
x>-1,\\
x\ne0,\\
x>-\dfrac23,\\
x(3x+1)<0.
\end{cases}
\]

Решаем последнее неравенство:

\[
x(3x+1)<0.
\]

Корни:

\[
x=0,
\qquad
x=-\frac13.
\]

\[
x\in\left(-\frac13;0\right).
\]

С учётом ограничений этот промежуток подходит полностью.

Ответ:

\[
\boxed{x\in\left(-\frac13;0\right)}
\]

\section*{11. Метод замены множителей в дробях}

Если логарифм стоит в числителе дроби, а справа ноль, метод можно применить к числителю.

Например:

\[
\frac{\log_x(5x-1)}{x^2-x-6}>0.
\]

Так как

\[
\log_x(5x-1)=\log_x(5x-1)-\log_x1,
\]

то логарифмическая часть заменяется на:

\[
(x-1)\bigl((5x-1)-1\bigr)
=
(x-1)(5x-2).
\]

Не забываем про ОДЗ:

\[
x>0,
\qquad
x\ne1,
\qquad
5x-1>0,
\qquad
x^2-x-6\ne0.
\]

А дальше решаем уже рациональное неравенство:

\[
\frac{(x-1)(5x-2)}{x^2-x-6}>0
\]

с учётом всех ограничений.

\section*{12. Типичные ошибки}

\begin{itemize}
  \item Вынести чётную степень из логарифма и забыть модуль.
  \item Забыть, что \((x-1)^2>0\) даёт \(x\ne1\), а не все \(x\).
  \item При корне чётной кратности поменять знак на интервалах.
  \item Записать два одинаковых корня как разные.
  \item Считать дискриминант там, где проще решить через неполное квадратное уравнение.
  \item Раскрывать модуль по определению вместо быстрого равносильного перехода.
  \item Перепутать систему и совокупность.
  \item Забыть ОДЗ логарифма с переменным основанием.
  \item В методе замены множителей забыть множитель \(a-1\).
  \item Представить \(0\) или \(1\) как логарифм с неправильным основанием.
  \item Нарисовать ось \(x\), хотя после замены переменной нужно рисовать ось новой переменной.
\end{itemize}

\section*{13. Короткая памятка}

\begin{center}
\renewcommand{\arraystretch}{1.35}
\begin{tabularx}{\linewidth}{@{}p{0.38\linewidth}X@{}}
\toprule
\textbf{Ситуация} & \textbf{Что делать} \\
\midrule
\(\log_a(f(x))^{2n}\) & Писать \(2n\log_a|f(x)|\), не забывая ОДЗ. \\
\((x-c)^2>0\) & Подходят все \(x\), кроме \(x=c\). \\
Корень чётной кратности & Знак на интервалах не меняется. \\
\(|f(x)|\le c\) & Переходить к двойному неравенству \(-c\le f(x)\le c\). \\
\(|f(x)|\ge c\) & Переходить к совокупности \(f(x)\ge c\) или \(f(x)\le -c\). \\
\(\log_a f\;\vee\;\log_a g\) & Использовать \((a-1)(f-g)\;\vee\;0\) плюс ОДЗ. \\
Справа \(1\) & Представить \(1=\log_a a\). \\
Справа \(0\) & Представить \(0=\log_a1\). \\
\bottomrule
\end{tabularx}
\end{center}

\newpage

\section*{Тренировочный блок}

Решите неравенства. Ответ запишите промежутками.

\subsection*{Средний уровень}

\begin{enumerate}
  \item
  \[
  \log_3(x-2)^2\le2.
  \]

  \item
  \[
  \log_2(x+1)^4>8.
  \]

  \item
  \[
  \log_5|x-3|\le1.
  \]
\end{enumerate}

\subsection*{Выше среднего}

\begin{enumerate}[resume]
  \item
  \[
  \log_{x}(2x-1)\ge1.
  \]

  \item
  \[
  \log_{x+1}(3x+2)<0.
  \]

  \item
  \[
  \log_{x-2}(5x-7)\le1.
  \]

  \item
  \[
  \frac{\log_x(4x-1)}{x-3}>0.
  \]

  \item
  \[
  \frac{\log_{x+2}(2x+5)}{x-1}\ge0.
  \]

  \item
  \[
  \log_3(x-1)^{10}\ge20.
  \]
\end{enumerate}

\subsection*{Сложная задача}

\begin{enumerate}[resume]
  \item
  \[
  \frac{\log_x(5x-1)\cdot\log_{2x-3}(x+4)}{x^2-5x+6}>0.
  \]
\end{enumerate}

\newpage

\section*{Ответы к тренировочному блоку}

\begin{center}
\renewcommand{\arraystretch}{1.45}
\begin{tabularx}{\linewidth}{@{}cX@{}}
\toprule
\textbf{№} & \textbf{Ответ} \\
\midrule
1 & \(x\in[-1;2)\cup(2;5]\) \\
2 & \(x\in(-\infty;-17)\cup(15;+\infty)\) \\
3 & \(x\in[-2;3)\cup(3;8]\) \\
4 & \(x\in\left[\dfrac12;1\right)\cup[2;+\infty)\) \\
5 & \(x\in\left(-\dfrac13;0\right)\) \\
6 & \(x\in\left(\dfrac75;2\right)\cup[3;+\infty)\) \\
7 & \(x\in\left(\dfrac14;1\right)\cup(3;+\infty)\) \\
8 & \(x\in\left[-\dfrac32;-1\right)\cup(1;+\infty)\) \\
9 & \(x\in(-\infty;-8]\cup[10;+\infty)\) \\
10 & \(x\in\left(\dfrac32;2\right)\cup\left(3;+\infty\right)\) \\
\bottomrule
\end{tabularx}
\end{center}

\end{document}